\documentclass[12pt, a4paper]{article}

% ==========================================
% PREÂMBULO
% ==========================================
\usepackage[utf8]{inputenc}
\usepackage[T1]{fontenc}
\usepackage[brazilian]{babel}
\usepackage{geometry}
\geometry{a4paper, left=3cm, top=3cm, right=2cm, bottom=2cm}
\usepackage{setspace}
\onehalfspacing
\usepackage[font=small, labelfont=bf, justification=justified, format=hang]{caption}
\usepackage{graphicx}
\usepackage{amsmath}
\usepackage{amssymb}
\usepackage{booktabs}      % Tabelas profissionais (\toprule, \midrule, \bottomrule)
\usepackage{multirow}      % Mesclar linhas em tabelas
\usepackage{array}         % Controle de colunas em tabelas
\usepackage{xcolor}        % Cores no texto e tabelas
\usepackage{hyperref}      % Links e referências clicáveis
\hypersetup{
    colorlinks=true,
    linkcolor=black,
    citecolor=black,
    urlcolor=blue
}
\usepackage{listings}      % Blocos de código
\usepackage[expansion=false]{microtype}     % Melhor justificação do texto (desabilitado para evitar erro de expansão de fonte)
\usepackage[alf]{abntex2cite}
\usepackage{float} % Permite forçar a posição exata de tabelas e imagens com [H]
\usepackage{tikz}
\usetikzlibrary{arrows.meta, shapes.geometric, positioning, fit, backgrounds}

\begin{document}

% ==========================================
% RESUMO
% ==========================================
\newpage
\begin{abstract}
\noindent
A avaliação \textit{in silico} da toxicidade aguda sistêmica é fundamental para reduzir o atrito clínico e o uso empírico de animais no desenvolvimento de fármacos. O principal desafio metodológico reside na modelagem de um espaço químico heterogêneo aliado à esparsidade empírica da Dose Letal Mediana ($DL_{50}$). Este trabalho descreve o desenvolvimento do ToxiciTOOL, um \textit{pipeline} computacional baseado em Redes Neurais Profundas e Aprendizado Multi-Tarefa (\textit{Multi-Task Learning} - MTL) para a predição simultânea da $DL_{50}$ em seis combinações de espécie (rato e camundongo) e via de administração (oral, intravenosa e intraperitoneal). Foram processados dados toxicológicos do banco de dados TOXRIC, culminando em uma matriz de 70.373 compostos únicos. A representação molecular otimizada empregou \textit{Morgan Count Fingerprints} concatenados a descritores físico-químicos (ADME). Para anular o vazamento de dados estruturais (\textit{data leakage}), implementou-se estratégias rigorosas de particionamento, incluindo o \textit{Scaffold Split} e uma adaptação escalável do algoritmo de Butina (\textit{Butina-batch}). A arquitetura MTL incorporou uma \textit{Masked Huber Loss} ponderada por incerteza homocedástica, equilibrando assimetrias nos conjuntos de treinamento. Sob o crivo prospectivo realista do \textit{split} Butina, o modelo obteve Coeficiente de Determinação ($R^2$) médio global de 0,481 e Erro Absoluto Médio (MAE) de 0,363 em log-escala, com picos de predição ($R^2 = 0,590$) para vias sub-representadas, como a intravenosa em ratos. Os resultados demonstram que a arquitetura Multi-Tarefa generaliza padrões estruturais de vias ricas em dados para tarefas com severa limitação amostral, consolidando o ToxiciTOOL como uma ferramenta viável de \textit{screening} alinhada aos princípios éticos dos 3Rs.

O repositório de código-fonte e artefatos do projeto está disponível em: \url{https://doi.org/10.5281/zenodo.19672989}

\vspace{1em}
\noindent\textbf{Palavras-chave:} Toxicidade aguda. $DL_{50}$. \textit{Multi-Task Learning}. Redes neurais. QSAR. Fingerprints moleculares.
\end{abstract}

\end{document}